\documentclass[11pt]{article}
\usepackage[T1]{fontenc}
\usepackage{lmodern,amsmath,amssymb,amsthm}
\usepackage[a4paper,margin=27mm]{geometry}
\usepackage[hidelinks]{hyperref}
\newtheorem{proposition}{Proposition}
\newtheorem{corollary}[proposition]{Corollary}
\newcommand{\Int}{\operatorname{int}}
\newcommand{\conv}{\operatorname{conv}}
\newcommand{\Rot}{\operatorname{Rot}}
\title{Boundary multiplicities and mirror capacity in rotational Stein tilings}
\author{MATH-004 --- LLM-generated research draft}
\date{21 October 2026}
\begin{document}
\maketitle

\noindent\textbf{Status and dependency boundary.}
These are proposed deductions, not independently accepted theorems or enabled
software filters. The owner chose generalization instead of reviewing the
earlier MATH-002/003 proofs. We do not use the two-orbit exclusion as a premise.
Statements requiring connected outer traces explicitly assume \textup{(CT)}
below. Removing that premise in the Stein setting relies on the unaccepted
MATH-002 proof, not on an audit performed here. No novelty claim is made.

\section{Setting and the precise generalization question}

Let $D=\{x\in\mathbb R^2:|x|\leq1\}$ and $O=0$. Finitely many congruent
compact Jordan regions tile $D$, with pairwise disjoint interiors and $N>1$.
One tile $C$ has $O\in\Int C$. An actual rotation
$R=\Rot_O(2\pi/m)$, $m\geq2$, permutes the tiles. Reflections are allowed;
no prescribed graph, connected internal interfaces, boundary smoothness, or
maximality of $m$ is assumed. This is a complete disk, not a free patch.

The center tile is fixed by $R$ and every other tile omits $O$. Brouwer's fixed
point theorem applied to any noncentral Jordan tile fixed by a nonidentity
power of $R$ shows that all noncentral orbits have size $m$. Write
$N=1+km$. Intrinsic order-$m$ symmetry of every tile follows by conjugating
$R|_C$ with a congruence, including a reflected congruence.

Put $r=\max_{x\in C}|x|$. A farthest-point rotation orbit has mean zero, so
any enclosing disk of $C$ with center $a$ and radius $t$ satisfies
$t^2\geq r^2+|a|^2$. Thus $O$ is its circumcenter and all tiles have common
circumradius $r$ and an interior circumcenter. Moreover $r<1$: if $|u|=1$
for $u\in C$, any congruence $f(x)=Qx+a$ from $C$ onto another tile gives
\[
 1\geq\frac1m\sum_{j=0}^{m-1}|QR^ju+a|^2=1+|a|^2,
\]
so $a=0$, contradicting disjoint interiors. In particular
$P(\conv C)\leq2\pi r<2\pi$, where $P$ denotes convex perimeter.

For the boundary-counting statements, assume explicitly:
\begin{quote}
\textup{(CT)} Every nonempty tile intersection with $\partial D$ is a
closed connected circular arc, possibly a singleton.
\end{quote}
The replacement argument in \texttt{two\_orbit\_stein\_obstruction.tex},
Lemma \texttt{lem:circumtrace}, proposes to deduce (CT) in this setting;
it remains unaccepted and is not revalidated by this note.

An outer tile has a positive-length outer trace. Let $b\leq k$ be the number
of its rotational orbits. Under (CT), their traces partition the circle up to
finitely many endpoints. Indeed two Jordan tiles cannot both occupy the disk
side of an open boundary subarc. Their total trace length is $2\pi$.

\section{Weighted repetitions for arbitrarily many orbits}

Choose one representative $T_i$ and congruence $f_i:C\to T_i$ for each
outer orbit. Use $f_{ij}=R^jf_iR^{-j}$ for the other copies, and put
\[
 q_i=f_i^{-1}(O)\ne O,\qquad
 A_i=f_i^{-1}(T_i\cap\partial D)
     =C\cap\partial\overline B(q_i,1).
\]
We have $C\subset\overline B(q_i,1)$, and the other transported arcs in
this orbit are $R^j A_i$ with supporting centers $R^j q_i$.

Partition these $b$ representatives by $q_i\sim q_j$ if
$q_i=R^\ell q_j$ for some integer $\ell$. Denote the classes by $a$, their
sizes by $s_a$, and their common arc lengths by $\alpha_a>0$.
This is a chart-dependent sufficient matching partition; it is not claimed
to identify every centered congruence when the prototype has extra symmetries.
Each class nevertheless contains genuinely centered-congruent outer orbits:
if $q_j=R^\ell q_i$, then $f_jR^\ell f_i^{-1}$ fixes $O$ and maps
$T_i$ onto $T_j$.

\begin{proposition}[Weighted boundary multiplicity, conditional on (CT)]
\label{prop:multiplicity}
With these definitions,
\[
 m\sum_a s_a\alpha_a=2\pi,\qquad
 m\sum_a\alpha_a\leq P(\conv C)\leq2\pi r.
\]
Consequently the arc-length-weighted mean multiplicity satisfies
\[
 \frac{\sum_a s_a\alpha_a}{\sum_a\alpha_a}\geq\frac1r,
 \qquad
 \max_a s_a\geq\left\lceil\frac1r\right\rceil,
 \qquad b\geq\left\lceil\frac1r\right\rceil.
\]
The ceiling inequalities are non-strict at integer thresholds.
\end{proposition}

\begin{proof}
Circle coverage proves the first equality. Every supporting disk above
contains $C$, so its contact arc is on $\partial\conv C$. Equal supporting
disks give the same complete arc. Distinct supporting circles cannot intersect
in the relative interior of either arc: at a transverse crossing that arc
would leave the other containing disk; distinct equal-radius tangent disks
have only one common point, inconsistent with the interior of $C$.
Thus different arcs meet only at endpoints.

For each class there are exactly $m$ distinct supporting centers, since
$q_i\ne O$. Different classes give disjoint sets of centers. The union of
all distinct transported arcs therefore has length $m\sum_a\alpha_a$,
bounded by convex perimeter. Dividing the two sums proves the weighted bound;
an integer maximum is at least its ceiling. Also $\sum_a s_a=b$.
\end{proof}

For example, $r<1/2$ forces at least three outer orbits in a single such
matching class, not merely three unrelated outer orbits. This is a conditional
geometric requirement, not an immediately detectable contact-graph filter.

\section{A local mirror constraint that needs no trace premise}

Suppose an isometry $s$ fixes $O$, reverses orientation, and maps a noncentral
tile $A$ to a tile $B$ in a different $R$-orbit. Then $s$ is reflection in a
line through $O$. No assertion that $s$ preserves the whole tiling is needed.

\begin{proposition}[Mirror capacity, without (CT)]\label{prop:mirror}
For $m\geq3$, put
\[
 c_m=\cos(\pi/m),\qquad
 \lambda_m=\frac{c_m}{\sin(\pi/(2m))},\qquad
 F_m=1+\lambda_m^2+2\lambda_m c_m.
\]
The existence of the reflected pair above forces
\[
 r\leq F_m^{-1/2},\qquad N\geq F_m,
 \qquad k\geq\left\lceil\frac{F_m-1}{m}\right\rceil.
\]
These necessary conditions allow additional interior and outer orbits, arbitrary
internal contacts, and positive-area Jordan boundaries.
\end{proposition}

\begin{proof}
The $m$ reflections $R^js$ send $A$ to distinct tiles $R^jB$. None of their
axes can meet $\Int A$, since a fixed point would then belong to two tile
interiors. Connectedness confines $\Int A$ to one component of the complement
of these axes. Taking closures, $A$ lies in a closed wedge $W$ of angle
$\pi/m$. This conclusion concerns $A$ only, not the entire tiling.

Let $P\in\Int A$ be its circumcenter and $d=|P|$. The convex hull
$K=\conv A$ contains the regular $m$-gon formed by the intrinsic rotation
orbit of any farthest point at distance $r$ from $P$. Hence
\[
 \overline B(P,r c_m)\subset K\subset W.
\]
Set $\beta=\pi/(2m)$ and use wedge-bisector coordinates $P=(x,y)$.
The two side-line distances give
\[
 x\sin\beta-|y|\cos\beta\geq r c_m,
 \qquad d\geq x\geq\lambda_m r.
\]
Among the $m$ farthest-orbit directions, one makes angle at most $\pi/m$
with the direction of $P$. Since its point is in $A\subset D$,
\[
 1\geq d^2+r^2+2dr c_m\geq F_m r^2.
\]
This proves the circumradius bound. Each congruent tile is contained in a
disk of area $\pi r^2$. Measure subadditivity, not equal tile-area sharing,
gives $\pi\leq N\pi r^2$, hence $N\geq F_m$. Substituting $N=1+km$
gives the stated orbit bound.
\end{proof}

\begin{corollary}[Boundary amplification, conditional on (CT)]
\label{cor:amplification}
If such a mirror pair exists, some transported-arc class satisfies
\[
 s_a\geq\left\lceil\sqrt{F_m}\right\rceil,
 \qquad b\geq\left\lceil\sqrt{F_m}\right\rceil,
 \qquad k\geq\left\lceil\sqrt{F_m}\right\rceil.
\]
The large matching class need not contain the original reflected pair.
As $m\to\infty$, $\sqrt{F_m}/m\to2/\pi$ and
$F_m/m^2\to4/\pi^2$.
\end{corollary}

\begin{proof}
Combine Propositions \ref{prop:multiplicity} and \ref{prop:mirror}.
The limits follow from $m\sin(\pi/(2m))\to\pi/2$ and
$\cos(\pi/m)\to1$.
\end{proof}

\section{A rotational outer shell cannot be promoted silently}

\begin{proposition}[A pure rotational shell needs additional interior tiles]
\label{prop:shell}
Assume (CT). It is impossible for every noncentral tile to be outer and
for all these tiles to be rotations about $O$ of one fixed tile.
\end{proposition}

\begin{proof}
Let $M=N-1$. All $M$ traces have the same positive length, hence length
$2\pi/M$, and partition the circle into equal successive arcs. The unique
rotation $\rho$ by $2\pi/M$ mapping one such arc to the next maps the
corresponding tiles: each is a centered rotation of the fixed tile, and a
proper positive circular arc determines that rotation uniquely.
Thus $\rho$ permutes all noncentral tiles. Their union is invariant; the
closure of its complement in $D$ is $C$, so $\rho(C)=C$ as well.

The whole tiling now has order-$M$ rotational symmetry with just one
noncentral orbit. Proposition \ref{prop:multiplicity}, applied to this
larger rotation, would require $1\geq1/r$, contrary to $r<1$.
\end{proof}

With additional interior tiles only their union, not their individual
partition, is forced to inherit this outer-shell rotation. In particular,
two rotationally matched outer orbits can create an order-$2m$ outer shell
without proving that the center tile or the inner tile partition has that
symmetry. That missing implication is not supplied by this note.

\section{Scope, controls and next mathematical decision}

The mirror capacity applies even if the reflected pair is not outer. It uses
no connected-trace lemma, no sector-endpoint coincidence and no two-orbit
exclusion. The stronger matching-class count and pure-shell exclusion retain
the explicit (CT) dependency. None of these implies that every forced centered
matching is reflective, or rules out all rotational Stein tilings.

For $m=3$, an equilateral triangle placed with one vertex at $O$ and opposite
edge on the outer chord of a $\pi/3$ wedge has circumradius $1/\sqrt3$ and
attains the mirror-capacity geometry. It is not a Stein outer tile, since it
contains $O$; it checks the constants, not the existence of a tiling. For
$m=2$, $c_m=0$ and the inscribed-polygon disk degenerates: no useful bound is
claimed here in that case. The earlier separate lens draft is not used.

The next bounded question is the orientation-preserving branch with three
moving orbits. In particular, for $m=5$ the mirror branch would require at
least four outer orbits by Corollary \ref{cor:amplification}; a hypothetical
three-orbit candidate would therefore have only orientation-preserving
centered matchings. Here the threshold does not depend on rounding:
$F_5=(13+5\sqrt5)/2>9$, hence $\sqrt{F_5}>3$.
Surviving possibilities include two outer orbits and one
interior orbit, or three outer orbits not all in one direct matching class.
Neither is excluded by counting alone. Seek the missing geometric constraint
there rather than assuming the outer-shell symmetry extends to the interior.

Independent review and a focused novelty check remain separate pending work.
No experimental solver behavior is used to infer any of these statements.
\end{document}
